\centerline{\bf \Huge Матанализ 2}
\centerline{\bf \Huge Непрерывность}

\begin{flushright}
        \scriptsize{Если ваш матанализ так хорош, то почему до сих пор нет матанализ 2???}
\end{flushright}

\problem{1}
(a) Если функции $f$ и $g$ непрерывны в точке $a$, то $f+g$ и $f \cdot g$ также непрерывны в точке $a$. Более того, если $g(a) \neq 0$, то $f / g$ также непрерывна в точке $a$.

(b) Если функция $g$ непрерывна в точке $a$, а функция $f$ непрерывна в точке $g(a)$, то $f(g)$ непрерывна в точке $a$.

\problem{2}
(a) \textbf{Теорема о двух милиционерах.} Последовательности $a_n, b_n$ и $c_n$ таковы, что $a_n \leqslant b_n \leqslant c_n$. Оказалось, что $a_n$ и $c_n$ стремятся к одному и тому же числу $d$. Докажите, что $b_n$ стремится к $d$.

(b) Докажите, что монотонная последовательность сходится тогда и только тогда, когда какая-нибудь ее подпоследовательность сходится.

\problem{3}
(a) Пусть $f:[0,1] \rightarrow[0,1]$ непрерывная функция. Докажите, что найдется $x \in[0,1]$ такой, что $f(x)=x$.

(b) Пусть $f:[0,1] \rightarrow \mathbb{R}$ непрерывная функция, принимающая все значения из отрезка $[0,1]$.. Докажите, что найдется $x \in[0,1]$ такой, что $f(x)=x$. 

\textit{Примечание:} нельзя пользоваться \href{https://ru.wikipedia.org/wiki/Теорема_Брауэра_о_неподвижной_точке}{теоремой Брауэра} о неподвижной точке (а я знаю, что всем хотелось!)

\problem{4}
Докажите, что

$$
\lim _{n \rightarrow \infty} \sqrt[n]{n}=1
$$

\problem{5}
Найдите все непрерывные функции $f: \mathbb{R} \rightarrow \mathbb{R}$ такие, что для любых действительных $x$ и $y$, если разность $x-y$ рациональна, то и разность $f(x)-f(y)$ рациональна.